\documentclass[11pt]{article}
\usepackage[margin=0.9in]{geometry}
\usepackage{amsmath,amssymb,amsthm,booktabs}
\usepackage[T1]{fontenc}
\usepackage{lmodern,microtype}
\usepackage[colorlinks=true,linkcolor=blue,urlcolor=blue]{hyperref}
\newtheorem{theorem}{Theorem}[section]
\newtheorem{lemma}[theorem]{Lemma}
\newtheorem{proposition}[theorem]{Proposition}
\newtheorem{corollary}[theorem]{Corollary}
\theoremstyle{remark}
\newtheorem{remark}[theorem]{Remark}
\title{Toward a four-dimensional Yang--Mills construction\\
\large Notebook 39: normalized local replacement in the bare Wilson measure}
\author{Research proof draft}
\date{17 October 2026}
\begin{document}
\maketitle
\begin{abstract}
We complete a normalized replacement operation for arbitrary finite
collections of separated fixed boxes in the positive bare Wilson
measure. On a good boundary the exact conditional law is restricted
to its local gauge/minimizer window and divided by its actual
conditional probability. On a bad boundary the law is unchanged.
The operation preserves the exterior marginal, gauge invariance,
and total raw partition mass. A positive expansion accounts for
every discarded interior configuration.

An explicit coupling has discrepancy marks dominated by independent
Bernoulli variables of parameter \(a=Ce^{-cT^2}\). For the exterior
jointly with any selected set \(A\) of box interiors, the exact
total variation distance is
\(\mathbb E[1-\prod_{i\in A}(1-p_i)]\le a|A|\).
We also prove local entropy, sparse-defect, cluster,
bounded-source and connected-cumulant comparison bounds. All are
independent of the number of other boxes.

The probability inputs are the full normalized Wilson bounds of
Notebook 38, not Gaussian surrogates. Geometry and box size remain
fixed. The operation has not been identified with a nonlinear
coarse-average fiber or propagated through general mixed
renormalization histories. This completes a restricted
finite-measure theorem and gives a definite stopping point for
handoff; it does not complete continuum Yang--Mills existence.
\end{abstract}

\section*{Inputs and proof boundary}

The two quantitative inputs are the fixed-box conditional tail
theorem and joint boundary-plaquette probability theorem
proved in Notebook 38. Both concern the actual compact positive
Wilson measure and include their normalization. Their box
geometry, boundary gauge, adjoint-invariant window, period
conditions and constants are restated below.
The new argument is an exact normalized operation built from
those inputs, together with its local comparison estimates.
It does not use the earlier homogeneous determinant calculation.

The group may be \(SU(N)\) for these bare-measure statements,
with \(\tau=\operatorname{Tr}/N\). This does not extend the
earlier \(SU(2)\) multiscale construction to every group.
An internal audit is included in the research archive.
Such review and successful typesetting do not certify
all upstream arguments independently.

